% =====================================================
% 题型：立体几何单选题 (q_solid_pyramid_skew_angle_two.tex)
% =====================================================
% =====================================================
% 难度：★★ (中档)
% =====================================================
\newcommand{\qSolidPyramidSkewAngleTwoStem}{%
  已知正三棱锥 $A - BCD$，$AB = 3$，$BC = 2$，$F$ 为 $CD$ 中点，则异面直线 $AF$ 与 $BD$ 所成角的余弦值为%
}

\newcommand{\qSolidPyramidSkewAngleTwoOptions}{%
  \mcoptions[4]{$\dfrac{\sqrt{2}}{2}$}{$\dfrac{\sqrt{3}}{3}$}{$\dfrac{\sqrt{2}}{4}$}{$\dfrac{\sqrt{2}}{8}$}%
}

\newcommand{\qSolidPyramidSkewAngleTwoFigure}[1][1.0]{%
  \begin{tikzpicture}[scale=#1, >=stealth]
    % Coordinates for regular triangular pyramid A-BCD
    \coordinate (B) at (0,0);
    \coordinate (C) at (3.0,0);
    \coordinate (D) at (1.2,1.2);
    
    % Base center projection is not strictly needed for drawing if we just place A
    \coordinate (A) at (1.5, 3.5);
    
    % Midpoint F of CD
    \coordinate (F) at ($(C)!0.5!(D)$);
    
    % Visible edges
    \draw[thick] (A) -- (B) -- (C) -- cycle;
    
    % Hidden edges
    \draw[thick, dashed] (A) -- (D);
    \draw[thick, dashed] (B) -- (D);
    \draw[thick, dashed] (C) -- (D);
    \draw[thick, dashed] (A) -- (F); % AF
    
    % Points
    \foreach \p in {A,B,C,D,F} {
      \fill (\p) circle (1.2pt);
    }
    
    % Labels
    \node[above] at (A) {\small $A$};
    \node[below left] at (B) {\small $B$};
    \node[below right] at (C) {\small $C$};
    \node[above left] at (D) {\small $D$};
    \node[right] at (F) {\small $F$};
    
    \ifteacher
      % Emphasize basis vectors DA, DB, DC
      \draw[->, gray, thick, dashed] (D) -- ($(D)!0.6!(A)$);
      \draw[->, gray, thick, dashed] (D) -- ($(D)!0.6!(B)$);
      \draw[->, gray, thick, dashed] (D) -- ($(D)!0.6!(C)$);
    \fi
  \end{tikzpicture}%
}

\newcommand{\qSolidPyramidSkewAngleTwoAnswer}{D}

\newcommand{\qSolidPyramidSkewAngleTwoSolution}{%
  本题由于正三棱锥底面为正三角形，侧棱相等，不便直接建立空间直角坐标系，可采用基底向量法进行计算。
  
  以 $\overrightarrow{DA}$，$\overrightarrow{DB}$，$\overrightarrow{DC}$ 为空间向量的一组基底。
  
  因为正三棱锥 $A-BCD$ 的侧棱长为 $3$，底面边长为 $2$，所以：
  \[ |\overrightarrow{DA}| = 3,\quad |\overrightarrow{DB}| = 2,\quad |\overrightarrow{DC}| = 2. \]
  
  因为底面 $\triangle BCD$ 是正三角形，所以：
  \[ \langle\overrightarrow{DB}, \overrightarrow{DC}\rangle = 60^\circ \implies \overrightarrow{DB} \cdot \overrightarrow{DC} = 2 \times 2 \times \cos 60^\circ = 2. \]
  
  由于侧面 $\triangle ABD$ 和 $\triangle ACD$ 是等腰三角形，且腰长为 $3$，底边长为 $2$，由余弦定理有：
  \[ \cos\angle ADB = \cos\angle ADC = \frac{3^2 + 2^2 - 3^2}{2 \times 3 \times 2} = \frac{4}{12} = \frac{1}{3}. \]
  
  所以，数量积为：
  \[ \overrightarrow{DA} \cdot \overrightarrow{DB} = |\overrightarrow{DA}| \cdot |\overrightarrow{DB}| \cos\angle ADB = 3 \times 2 \times \frac{1}{3} = 2. \]
  \[ \overrightarrow{DA} \cdot \overrightarrow{DC} = |\overrightarrow{DA}| \cdot |\overrightarrow{DC}| \cos\angle ADC = 3 \times 2 \times \frac{1}{3} = 2. \]
  
  因为 $F$ 是 $CD$ 的中点，所以 $\overrightarrow{DF} = \dfrac{1}{2}\overrightarrow{DC}$。
  
  则向量 $\overrightarrow{AF}$ 可表示为：
  \[ \overrightarrow{AF} = \overrightarrow{DF} - \overrightarrow{DA} = \frac{1}{2}\overrightarrow{DC} - \overrightarrow{DA}. \]
  
  计算异面直线两方向向量 $\overrightarrow{AF}$ 与 $\overrightarrow{DB}$ 的数量积与模长：
  \[ \overrightarrow{AF} \cdot \overrightarrow{DB} = \left(\frac{1}{2}\overrightarrow{DC} - \overrightarrow{DA}\right) \cdot \overrightarrow{DB} = \frac{1}{2}\overrightarrow{DC} \cdot \overrightarrow{DB} - \overrightarrow{DA} \cdot \overrightarrow{DB} = \frac{1}{2}(2) - 2 = -1. \]
  
  计算 $|\overrightarrow{AF}|$ 模长的平方：
  \[ 
  |\overrightarrow{AF}|^2 = \left(\frac{1}{2}\overrightarrow{DC} - \overrightarrow{DA}\right)^2 = \frac{1}{4}|\overrightarrow{DC}|^2 - \overrightarrow{DC} \cdot \overrightarrow{DA} + |\overrightarrow{DA}|^2 = \frac{1}{4}(4) - 2 + 9 = 8.
  \]
  所以 $|\overrightarrow{AF}| = 2\sqrt{2}$。而 $|\overrightarrow{DB}| = 2$。
  
  设异面直线 $AF$ 与 $BD$ 所成角为 $\theta$，则有：
  \[ \cos\theta = \frac{|\overrightarrow{AF} \cdot \overrightarrow{DB}|}{|\overrightarrow{AF}| \cdot |\overrightarrow{DB}|} = \frac{|-1|}{2\sqrt{2} \times 2} = \frac{\sqrt{2}}{8}. \]
  
  故选 \textbf{D}。
}

\newcommand{\qSolidPyramidSkewAngleTwoDiagnosis}{%
  用于课堂精准变式训练。正三棱锥是不对称且不垂直的几何模型，直接建系较为复杂。本题考查学生在非垂直环境下使用“基底向量法”处理数量积、夹角的能力。
}

\newcommand{\qSolidPyramidSkewAngleTwoTeachingNotes}{%
  \item \textbf{方法总结}：正三棱锥（或没有三条两两垂直棱的图形）中，首选以顶点或某个底面顶点为起点的三条棱作为“空间基底”解决夹角。
  \item \textbf{易错点}：基底间夹角的余弦值容易漏乘两向量模长，需板书步骤强调。
}
